Campaign 3 establishes what KV-cache geometry can and cannot detect through 28 experiments across three architectures.

\paragraph{What works.} Four detection capabilities are demonstrated, unified by a common mechanism of \emph{output suppression}: deception (AUROC $\approx 1.0$ within-model, $\approx 0.86$ cross-model), safety refusal (AUROC $\approx 0.91$), impossibility refusal (AUROC $\approx 0.94$), and absence-of-suppression in abliterated models (AUROC $\approx 0.88$ vs.\ benign). All involve states where the model withholds its full representational capacity. The key evidence is Experiment 36: impossibility refusal (AUROC $\approx 0.94$) is \emph{more} detectable than safety refusal (AUROC $\approx 0.91$), and the two refusal types are geometrically indistinguishable (AUROC $\approx 0.65$)---proving the signal is about withholding, not content type. Extended features (9 total) yield AUROC $\approx 0.95$, though the improvement over 4 features is within CV noise at this sample size \citep{varoquaux2018cross}. A fifth candidate---sycophancy detection---was refuted: the apparent signal (AUROC 0.9375) was a prompt-length confound, and the length-matched effect ($d = 0.107$) requires $n \approx 1{,}400$ per condition (\Cref{sec:sycophancy}).

\paragraph{What doesn't work.} Five independent experiments show that confabulation has no \emph{practically reliable} KV-cache signature. Aggregate features produce negligible effect sizes ($d = 0.052$). A directional signal exists in per-layer profiles (effective AUROC $= 0.977$, $p = 0.02$ by permutation test) but is too fragile for classification ($p = 0.21$ for logistic regression). An initially promising contrastive signal ($d = 1.081$) was debunked by a 4-condition control as an information volume artifact. The geometric signal is real but unstable: confabulation shares the misalignment axis direction but its perturbation drifts $41.2^{\circ}$ during generation, compared to $2.1^{\circ}$ for deception. The signal exists; it just cannot be reliably exploited.

\paragraph{What this means.} The output suppression framework provides a principled boundary for internal state monitoring. Any approach that reads internal representations---probing classifiers, representation engineering, sparse autoencoders, or KV-cache analysis---should be able to detect \emph{output suppression} (the model withholding its full representational capacity) but not \emph{ignorance} (the model lacking information). The confabulation signal is real at a geometric level but too unstable for deployment; practical confabulation detection requires external verification (retrieval-augmented generation, citation checking, knowledge graph alignment), not internal monitoring alone. This finding should inform how the AI safety community allocates research effort: internal monitoring for detecting when models are withholding, external verification for detecting when models are confidently wrong.

\paragraph{The geometry is universal.} KV-cache category geometry is preserved from 0.6B to 70B parameters ($\rho = 0.83$--$0.90$), invariant to GPU hardware ($r > 0.999$), and transferable across architectures via linear classification (AUROC $\approx 0.86$). A single ``suppression axis'' captures deception within $8.4^{\circ}$ of sycophancy and confabulation directions; however, only active suppression states produce stable perturbations along this axis. Confabulation's perturbation drifts $41.2^{\circ}$ during generation, and the sycophancy signal was confounded by prompt length (\Cref{sec:sycophancy}). Honest processing is geometrically richer than suppressive processing by four independent measures. These are structural properties of how transformers organize information in the KV-cache, not artifacts of specific models or hardware.

\paragraph{Open problems.} Cross-model transfer, while solved for deception (AUROC $\approx 0.86$), remains untested for refusal and output suppression detection; the minimum observed cross-model AUROC is 0.14, indicating that transfer is not guaranteed across all architecture pairs. Sycophancy detection via aggregate KV-cache features remains an open problem; the length-matched effect ($d = 0.107$) requires either per-layer analysis from late layers, substantially larger prompt batteries ($n \approx 1{,}400$), or both \citep{vennemeyer2025sycophancy}. The SAE-based comparison (sparse autoencoder feature activation versus aggregate cache geometry) is deferred to future work. And the fundamental question---whether naturally emergent deception (not instructed via system prompt) produces the same output suppression geometry as instructed deception---awaits models that demonstrably deceive spontaneously.
